% Generated draft: requires mathematical and source review.
\documentclass[11pt]{article}
\usepackage[margin=25mm]{geometry}
\usepackage{fontspec}
\setmainfont{Latin Modern Roman}
\usepackage{amsmath,amssymb,mathtools,unicode-math}
\setmathfont{Latin Modern Math}
\usepackage{enumitem,needspace,xurl,hyperref}
\setlength{\parindent}{0pt}
\setlength{\parskip}{5pt}
\setlength{\emergencystretch}{3em}
\newcommand{\N}{\mathbb{N}}
\newcommand{\Z}{\mathbb{Z}}
\newcommand{\Q}{\mathbb{Q}}
\newcommand{\R}{\mathbb{R}}
\newcommand{\C}{\mathbb{C}}
\newcommand{\F}{\mathbb{F}}
\newcommand{\A}{\mathbb{A}}
\newcommand{\PP}{\mathbb{P}}
\newcommand{\E}{\mathbb{E}}
\newcommand{\eps}{\varepsilon}
\newcommand{\dd}{\,\mathrm{d}}
\begin{document}
\section*{UnsolvedMath ID: 10400011}
\textbf{Source problem code:} catalogue entry
\subsection{Definitions and mathematical statement}
For an oriented link \(L\subset S^3\), Khovanov categorification assigns a bigraded homology, with homological grading \(i\) and polynomial grading \(j\), whose graded Euler characteristic recovers a chosen normalization of the Jones polynomial. Over \(\Q\), this relationship has the form
\[
 \sum_{i,j}(-1)^i q^j\dim_\Q\operatorname{Kh}^{i,j}(L;\Q)=J_L(q),
\]
where \(J_L\) denotes the matching Jones normalization and variable convention. The reduced and unreduced versions and grading shifts must be distinguished; the excerpt does not choose them.

Problem 1.11 asks to \emph{understand Khovanov's categorification of the Jones polynomial}. This is a qualitative request for mathematical understanding of an existing construction, not a stated existence conjecture for a bigraded theory. It gives no particular geometric interpretation, structural theorem, computation or criterion of understanding to prove. The Euler-characteristic display explains categorification but does not exhaust the requested understanding.

The registry labels the problem solved on the ground that Khovanov homology is established and has substantial applications. That report does not identify which more specific conceptual objective, if any, the source intended by ``understand''. The draft therefore retains the programme rather than silently replacing it with the existence of Khovanov homology or with a new prescribed theorem.
\subsection{Short English statement}
Seek mathematical understanding of Khovanov's homological refinement of the Jones polynomial. The source states a broad programme and does not specify a further theorem whose proof would settle it.
\subsection{Sources}
\begin{itemize}
\item Local UnsolvedMath registry, permanent ID 10400011, source code catalogue entry; no proposer or proposal year supplied; updated 2026-08-06; dated status solved.
\item Ohtsuki, \emph{Problems on invariants of knots and 3-manifolds} (2002), Problem 1.11, PDF page 18: \url{https://msp.org/gtm/2002/04/gtm-2002-04-024s.pdf}. Registry access date 2026-07-29; PDF font-structure extraction. Provenance distinguishes 231 numbered bold problem headings from surrounding references. Present primary-literature status and redistribution terms require review.
\item Khovanov is the supplied categorification attribution. The research summary mentions Khovanov homology, the Rasmussen invariant, concordance and four-dimensional topology without exact references or a specific resolution of the source's qualitative objective.
\end{itemize}
\end{document}
